% HTB: independent completed-Hurwitz / original-theta calculation.
% This file is a chapter source; the parent owns its cumulative wrapper.
\section{The exact completed Hurwitz--theta bridge over the full support lattice}
\label{sec:completed-hurwitz-theta-bridge}

\paragraph{Original-object correction, 23 September 2026.}
The auxiliary functions in this chapter are compared with the original
$\zeta$ and $F_t$; they are not replacements for them. HTB13 already
provides the fractional-module isomorphism at the exceptional points.
CFD1--10 retains its comparison with the original regular lattice and
all finite jets. ZTC13--16 proves the joint inverse with the boundary
component kept. Projection to $H_0$ alone discards the actual Hurwitz
deformation. The original arithmetic receiving formulas and both
Mellin endpoint terms are given in ZTC1--21.

\paragraph{Objects and source provenance.}
The shifted family below is the one in the supplied April 2026 programme
draft \emph{Split-Zero Support, Residue-Parity Reversal, Shifted-Hurwitz CUE
Defects, and Finite Phase-Space Enlargement}, printed byline ``Compiled
theorem draft'', source labels \texttt{thm:universal-flow} and
\texttt{thm:zero-motion}. The actual source is
\texttt{Chatnotes/globalization cue/}\allowbreak
\nolinkurl{split_zero_cue_shifted_hurwitz_phase_space.tex},
lines 602--630 and 1003--1060; the exact retained copy is under
\texttt{hurwitz\_intake/programme\_source/}.
No human author is specified in those supplied
bytes, and none is inferred here. Its retained source has SHA-256
\nolinkurl{52f2a6bccdb15213b1992abef9f8d6476d6e9ea5566ee7ed9d85c9a2d7cc2705}.
The heat family and every factor in its coordinate map are those of
Brad Rodgers and Terence Tao,
\href{https://arxiv.org/abs/1801.05914v5}{\emph{The de Bruijn--Newman
constant is non-negative}, arXiv:1801.05914v5}, original source
\nolinkurl{fmp-template.tex}, equations \texttt{hoz}, \texttt{sas},
\texttt{phidef}, and \texttt{htdef}, lines 78--101. That source credits
N. C. de Bruijn's \emph{The roots of trigonometric integrals} (1950),
C. M. Newman's \emph{Fourier transforms with only real zeroes} (1976),
and the zero dynamics of G. Csordas, W. Smith, and R. S. Varga (1994).
Those earlier papers have not been independently read for this chapter;
their attribution is inherited explicitly from the original Rodgers--Tao
source. The theta identity used here is proved below, including its
Mellin factors. The full support carrier is the one in
\href{https://github.com/KokunoYumeto/zeta-function-research-reader/blob/a2851f9eb352e7e4376bc629de86fa4ed9c1c434/workbenches/splitzero-tandem/foundations/split-support-geometry-v11/split_support_geometry_arithmetic_curve_v11.tex}{\emph{Split Support Geometry}, v11,
Definition \texttt{def:lattice-split}}. Its operations and all maps needed
here are stated and verified in the final subsection.

\subsection{Two parameters and the original coordinate map}

Use $t$ for the Hurwitz shift and $h$ for heat time. They are independent
parameters. Put
\begin{equation}\tag{HTB1}
\begin{gathered}
 a=1+t,\qquad \operatorname{Re}a>0,\qquad
 F_t(s)=\zeta(s,a),\\
 A(s)=\frac{s(s-1)}2\pi^{-s/2}\Gamma(s/2),\qquad
 \Xi_t(s)=A(s)F_t(s),\\
 s(z)=\frac12+\frac{iz}{2},\qquad
 z(s)=-2i\left(s-\frac12\right),\qquad
 H_t^{\mathrm{Hur}}(z)=\frac18\Xi_t(s(z)).
\end{gathered}
\end{equation}
For complex $a$ in the indicated half-plane, $(n+a)^{-s}$ uses the
holomorphic logarithm with argument in $(-\pi/2,\pi/2)$. All subsequent
identities in $s$ are identities of meromorphic functions; values at
removable singularities mean their analytic extensions. Statements about
real parameters explicitly use $t>-1$.

Here is a direct continuation of the Hurwitz input, so that neither its
pole nor the endpoint values are assumed in the comparison. Define
Bernoulli polynomials by
\begin{equation}\tag{HTB2}
 \frac{x e^{b x}}{e^x-1}
 =\sum_{k=0}^{\infty}B_k(b)\frac{x^k}{k!},\qquad |x|<2\pi.
\end{equation}
For $\operatorname{Re}s>1$, termwise integration gives
\begin{equation}\tag{HTB3}
 \Gamma(s)F_t(s)
 =\int_0^{\infty}x^{s-1}\frac{e^{-a x}}{1-e^{-x}}\,dx.
\end{equation}
The interchange is absolutely convergent on compact subsets of
$\operatorname{Re}s>1$, $\operatorname{Re}a>0$: near zero the absolute
integrand is bounded by a constant times $x^{\operatorname{Re}s-2}$,
and at infinity it decays exponentially. Subtracting the first $N+1$
terms of HTB2 on $(0,1)$ gives, for $\operatorname{Re}s>-N$,
\begin{equation}\tag{HTB4}
\begin{split}
 F_t(s)=\frac1{\Gamma(s)}\bigg\{
 &\int_1^{\infty}x^{s-1}\frac{e^{-a x}}{1-e^{-x}}\,dx\\
 +&\int_0^1x^{s-1}
 \left(\frac{e^{-a x}}{1-e^{-x}}
       -\sum_{k=0}^{N}\frac{B_k(1-a)}{k!}x^{k-1}\right)dx\\
 +&\sum_{k=0}^{N}\frac{B_k(1-a)}{k!(s+k-1)}\bigg\}.
\end{split}
\end{equation}
The remainder is $O(x^N)$ locally uniformly in $a$. Differentiating in
$a$ and $s$ preserves locally integrable majorants on smaller compact
sets. Thus HTB4 gives joint meromorphic continuation, holomorphic in
$a$, with compatible values as $N$ increases.

The recurrence $\Gamma(s+1)=s\Gamma(s)$ and $\Gamma(1)=1$ show that
$\Gamma$ has residue $(-1)^n/n!$ at $-n$. Its reciprocal is entire with
simple zeros there and no other zeros: for example this follows directly
from the locally uniformly convergent Euler product
$1/\Gamma(s)=s e^{\gamma s}\prod_{n\ge1}(1+s/n)e^{-s/n}$, whose
nonvanishing away from the listed factors follows from the absolute
convergence of the logarithmic tail. In HTB4 these reciprocal zeros
cancel every displayed pole except the one at $s=1$. At $s=-n$ only
the term $k=n+1$ contributes to the value. Replacing $x$ by $-x$ in
HTB2 proves $B_k(1-a)=(-1)^k B_k(a)$. Consequently
\begin{equation}\tag{HTB5}
 \operatorname*{Res}_{s=1}F_t(s)=1,\qquad
 F_t(-n)=-\frac{B_{n+1}(a)}{n+1}\quad(n\ge0),\qquad
 \partial_t F_t(s)=-sF_t(s+1).
\end{equation}
For the last equality differentiate the convergent series on
$\operatorname{Re}s>1$; HTB4 and uniqueness of meromorphic continuation
extend it everywhere. In particular $F_t(0)=1/2-a$.

\subsection{The theta integral with all original constants}

For $x>0$, write
\[
 \psi(x)=\sum_{n=1}^{\infty}e^{-\pi n^2x},\qquad
 \Theta(x)=1+2\psi(x).
\]
The Gaussian identity gives
\begin{equation}\tag{HTB6}
 \Theta(x)=x^{-1/2}\Theta(1/x),\qquad
 \psi(1/x)=\sqrt{x}\,\psi(x)+\frac{\sqrt{x}-1}{2}.
\end{equation}
One direct proof of the required Gaussian identity is to periodize
$e^{-\pi x y^2}$. The periodized series and all its derivatives converge
uniformly on $[0,1]$. Its $k$th Fourier coefficient is
$\int_{\mathbb R}e^{-\pi x y^2}e^{-2\pi i k y}dy
=x^{-1/2}e^{-\pi k^2/x}$. This integral follows by scaling from
$\int e^{-\pi y^2}dy=1$ and differentiating its Fourier transform:
integration by parts gives $\widehat f'(\eta)=-2\pi\eta\widehat f(\eta)$,
with $\widehat f(0)=1$. The absolutely convergent Fourier series of the
periodization evaluated at zero proves HTB6.

Define the original kernel, without changing its argument or amplitude,
by
\begin{equation}\tag{HTB7}
\begin{split}
 \Phi(u)
 &=\sum_{n=1}^{\infty}
   \left(2\pi^2n^4e^{9u}-3\pi n^2e^{5u}\right)
       e^{-\pi n^2e^{4u}}\\
 &=e^u\left(2x^2\psi''(x)+3x\psi'(x)\right)_{x=e^{4u}}.
\end{split}
\end{equation}
If $g(u)=e^u\psi(e^{4u})$, direct differentiation and HTB6 give
\begin{equation}\tag{HTB8}
 \Phi(u)=\frac18\bigl(g''(u)-g(u)\bigr),\qquad
 g(-u)=g(u)+\sinh u,\qquad \Phi(-u)=\Phi(u).
\end{equation}
Indeed $g''-g=24e^{5u}\psi'(e^{4u})+16e^{9u}\psi''(e^{4u})$,
and $\partial_u^2-1$ annihilates $\sinh u$ and commutes with reflection.
For $u\ge0$, each summand in HTB7 is strictly positive because
$2\pi n^2 e^{4u}-3>0$. Moreover the series and each fixed number of its
derivatives are bounded by a constant times $e^{C u}e^{-\pi e^{4u}}$:
factor out $e^{-\pi e^{4u}}$ and bound the remaining sums of powers of
$n$ times $e^{-\pi(n^2-1)}$. Evenness gives the corresponding estimate
as $u\to-\infty$. These bounds justify all differentiations and
integrations that follow.

For $\operatorname{Re}s>1$, substitute $x=e^{4u}$ and integrate by parts
twice. The endpoint terms vanish, at infinity exponentially and at zero
by HTB6 and $\operatorname{Re}s>1$. The exact calculation is
\begin{equation}\tag{HTB9}
\begin{split}
 \int_{\mathbb R}\Phi(u)e^{(2s-1)u}du
 &=\frac14\int_0^{\infty}
       \left(2x^{s/2+1}\psi''(x)+3x^{s/2}\psi'(x)\right)dx\\
 &=\frac{s(s-1)}8\int_0^{\infty}\psi(x)x^{s/2-1}dx\\
 &=\frac{s(s-1)}8\pi^{-s/2}\Gamma(s/2)\zeta(s)
 =\frac14\xi(s).
\end{split}
\end{equation}
For clarity, with $b=s/2$ the integration-by-parts multiplier is
$\frac14(2b(b+1)-3b)=s(s-1)/8$. At zero HTB6 gives
$\psi(x)=\frac12x^{-1/2}-\frac12+O(x^{-1/2}e^{-\pi/x})$, and the
derivatives give the claimed vanishing boundary terms. The last
interchange in HTB9 is absolutely convergent for $\operatorname{Re}s>1$.
The left-hand side is entire by the preceding super-exponential bounds.
It therefore supplies the entire continuation of $\xi=A\zeta$ and,
by evenness, proves $\xi(s)=\xi(1-s)$. All its coefficients respect
complex conjugation because the kernel is real.

The original heat family is thus
\begin{equation}\tag{HTB10}
\begin{gathered}
 H_h^{\mathrm{heat}}(z)=\int_0^{\infty}e^{h u^2}\Phi(u)\cos(zu)\,du,
 \qquad (h,z)\in\mathbb C^2,\\
 H_0^{\mathrm{heat}}(z)=\frac18\xi(s(z))=H_0^{\mathrm{Hur}}(z),\qquad
 \partial_h H_h^{\mathrm{heat}}=-\partial_z^2 H_h^{\mathrm{heat}}.
\end{gathered}
\end{equation}
Joint holomorphy and the differential equation follow from locally
uniform integrable majorants for every $h,z$ derivative. To obtain the
middle identity, put $2s-1=iz$ in HTB9 and use evenness: the integral
over the whole real axis is twice the cosine integral. Hence the factor
is exactly $1/8$. Every $H_h^{\mathrm{heat}}$ is even; for real $h$ it
also satisfies $H_h^{\mathrm{heat}}(\bar z)
=\overline{H_h^{\mathrm{heat}}(z)}$.

\subsection{An invertible meromorphic map and its exact product}

Let $\mathcal M_s$ and $\mathcal M_z$ be the fields of meromorphic
functions on the entire $s$-plane and $z$-plane, respectively. The
following maps are everywhere defined as maps of these fields viewed
as complex vector spaces:
\begin{equation}\tag{HTB11}
 U:\mathcal M_s\longrightarrow\mathcal M_z,\qquad
 (Uf)(z)=\frac{A(s(z))}{8}f(s(z)),\qquad
 (U^{-1}G)(s)=\frac{8G(z(s))}{A(s)}.
\end{equation}
Since $A$ is a nonzero meromorphic function and $s,z$ are inverse affine
coordinates, these formulas prove complex linearity and both inverse
identities. With ordinary pointwise products $U$ is not multiplicative:
its two product expressions differ by a factor $A/8$. The exact
transported product and unit on the same meromorphic vector space are
\begin{equation}\tag{HTB12}
 (G\mathbin{\odot}J)(z)=\frac{8G(z)J(z)}{A(s(z))},\qquad
 \mathbf 1_{\odot}(z)=\frac{A(s(z))}{8},\qquad
 U(fq)=Uf\mathbin{\odot}Uq.
\end{equation}
All field axioms follow either by this inverse transport or by the
displayed multiplication formula. Thus $U$ is a unital field
isomorphism to $(\mathcal M_z,\odot)$, as well as a linear isomorphism
to the ordinary meromorphic receiver. These two assertions use
different, explicitly stated products. In particular the ordinary
constant function $1$ is not asserted to be its transported unit.

There is an exact local regularity rule. At a finite $s_0$, put
$z_0=z(s_0)$ and let $\nu_{s_0}(f)$ denote meromorphic order, with
$\nu(0)=+\infty$. Since the coordinate derivative $ds/dz=i/2$ is
nonzero,
\begin{equation}\tag{HTB13}
 \nu_{z_0}(Uf)=\nu_{s_0}(A)+\nu_{s_0}(f),\qquad
 \mathcal D_{s_0}
 =\{f:\nu_{s_0}(f)\ge-\nu_{s_0}(A)\}
 \xrightarrow[\;U\;]{\ \cong\ }\mathcal O_{z_0}.
\end{equation}
Here the domain is a fractional module of meromorphic germs and the
arrow is a module isomorphism over the affine coordinate isomorphism;
it is not claimed to be a subring with the ordinary product. At a point
where $A$ is a holomorphic unit, it restricts to an isomorphism of
holomorphic modules and satisfies the exact evaluation formula
\[
 \operatorname{ev}_{z_0}Uf
   =\frac{A(s_0)}8\operatorname{ev}_{s_0}f.
\]
At a zero or pole of $A$, HTB13, rather than an undefined product of
separate point values, specifies the evaluation domain.

For example, if $A$ is a unit at $s_0$ and $F_t$ has a zero there,
ordinary affine substitution, denoted $C(f)=f\circ s$, induces the
unital algebra isomorphism
\begin{equation}\tag{HTB14}
 \mathcal O_{s_0}/(F_t)
 \xrightarrow[\;C\;]{\ \cong\ }
 \mathcal O_{z_0}/(H_t^{\mathrm{Hur}}).
\end{equation}
Indeed $U F_t=(A\circ s)C(F_t)/8$ and $A\circ s$ is a unit, so the
two target principal ideals coincide. This proves equality of local
multiplicities and the correspondence of every jet, without changing
their coordinate scale. Explicitly, for $w=z-z_0$,
\[
 (Uf)(z_0+w)=\frac18\sum_{n\ge0}
       \left(\frac{i}{2}\right)^n w^n
       \sum_{j=0}^n\frac{A^{(j)}(s_0)f^{(n-j)}(s_0)}{j!(n-j)!}.
\]
Thus the complete finite jet map is triangular, with nonzero diagonal
$A(s_0)(i/2)^n/8$. It is $C$, not multiplication by that unit, which
gives the ordinary unital algebra arrow in HTB14.

\subsection{Both infinitesimal generators and their difference}

Define $D:\mathcal M_s\to\mathcal M_s$ by $Df(s)=s f(s+1)$.
All differential and shift operators in the next display have the
entire meromorphic field as their algebraic domain; no bounded-operator
or semigroup assertion on an unspecified normed space is being used.
Direct use of HTB11 gives
\begin{equation}\tag{HTB15}
\begin{split}
 (D_H G)(z):=(UDU^{-1}G)(z)
 &=c(s(z))G(z-2i),\\
 c(s)=\frac{sA(s)}{A(s+1)}
 &=\sqrt{\pi}\,\frac{s(s-1)}{s+1}
       \frac{\Gamma(s/2)}{\Gamma((s+1)/2)},\\
 \partial_tH_t^{\mathrm{Hur}}&=-D_HH_t^{\mathrm{Hur}}.
\end{split}
\end{equation}
The sign and displacement follow from $s(z-2i)=s(z)+1$, and the last
identity follows from HTB5. Gamma recurrence gives the removable values
$c(0)=-2$, $c(-1)=-2\pi$, and $c(1)=0$. The complete product in HTB15
is meromorphic: at negative odd integers below $-1$ the vanishing
multiplier may multiply a pole of the shifted function, and evaluation
is performed after multiplication. For instance, at $s=0$ the formula
will give $\partial_tH_t^{\mathrm{Hur}}(i)=1/8$, exactly as HTB20 below.

The original heat generator $B=-\partial_z^2$ has pullback
\begin{equation}\tag{HTB16}
 \mathcal L_H:=U^{-1}BU,\qquad
 \mathcal L_H f(s)=\frac1{4A(s)}\frac{d^2}{ds^2}\bigl(A(s)f(s)\bigr).
\end{equation}
This includes both powers from $ds/dz=i/2$ and the minus sign in $B$;
the original factors $8$ in $U$ and $U^{-1}$ cancel each other.
The actual pulled-back heat trajectory is
$f_h(s)=8H_h^{\mathrm{heat}}(z(s))/A(s)$, with $f_0=\zeta$ and
$\partial_hf_h=\mathcal L_H f_h$. No assertion is needed that an
operator exponential converges on every meromorphic function.

The two actual generators in the same $s$-space are $\mathcal L_H$ and
$-D$. Their exact difference is the operator
\begin{equation}\tag{HTB17}
\begin{split}
 \mathcal E f(s)
 &:=(\mathcal L_H+D)f(s)\\
 &=\frac14 f''(s)+\frac{b(s)}2f'(s)
   +\frac{b'(s)+b(s)^2}{4}f(s)+s f(s+1),\\
 b(s)&=\frac{A'(s)}{A(s)}
 =\frac1s+\frac1{s-1}-\frac{\log\pi}{2}
       +\frac12\frac{\Gamma'(s/2)}{\Gamma(s/2)},\\
 (B+D_H)U&=U\mathcal E.
\end{split}
\end{equation}
The expansion follows by applying the ordinary product rule twice to
HTB16. Every displayed coefficient is interpreted meromorphically;
apparent singularities from the expanded expression do not change the
domain just specified. HTB17 is the exact map of the generator defect,
not a declaration that the two flows are unrelated.

\subsection{All poles, the entire parameters, and reflection}

The only finite singularities of $A$ are simple poles at $s=-2m$,
$m\ge1$. The apparent pole at zero is removable and $A(0)=-1$;
$s=1$ is its only zero, and
$A(s)=(s-1)/2+O((s-1)^2)$ there. These statements follow from the
gamma pole calculation following HTB4, the nonvanishing of gamma
away from its poles, and the two polynomial factors of $A$. Hence
HTB5 gives
\begin{equation}\tag{HTB18}
\begin{gathered}
 \Xi_t(0)=t+\frac12,\qquad \Xi_t(1)=\frac12,\\
 a_m:=\operatorname*{Res}_{s=-2m}A(s)
   =\frac{2(-1)^m(2m+1)\pi^m}{(m-1)!},\\
 r_m(t):=\operatorname*{Res}_{s=-2m}\Xi_t(s)
   =\frac{2(-1)^{m+1}\pi^m}{(m-1)!}B_{2m+1}(1+t).
\end{gathered}
\end{equation}
These are all possible poles of $\Xi_t$. A listed point is a pole
exactly when $r_m(t)\ne0$; otherwise it is removable and its value is
$a_m\partial_sF_t(-2m)$. To check every constant, the residue of
$\Gamma(s/2)$ in the $s$-coordinate is $2(-1)^m/m!$ and
$(-2m)(-2m-1)/2=m(2m+1)$. Multiplying their product by
$F_t(-2m)=-B_{2m+1}(1+t)/(2m+1)$ gives HTB18. The endpoint at one
uses the residue $1$ of $F_t$, not its undefined point value.

In particular the first residue and the exact classification are
\begin{equation}\tag{HTB19}
\begin{gathered}
 r_1(t)=2\pi(1+t)\left(t+\frac12\right)t,\\
 t\in\mathbb R,\ t>-1:\qquad
 \Xi_t\text{ is entire}\quad\Longleftrightarrow\quad
 t\in\left\{0,-\frac12\right\},\\
 \Xi_0(s)=\xi(s),\qquad
 \Xi_{-1/2}(s)=(2^s-1)\xi(s).
\end{gathered}
\end{equation}
Indeed $B_3(a)=a(a-1/2)(a-1)$ follows by expanding HTB2. Thus the
first pole alone excludes every other $t>-1$. At zero, entireness was
proved in HTB9. At $t=-1/2$, the absolutely convergent identity
$\zeta(s,1/2)=2^s\sum_{n\ge0}(2n+1)^{-s}=(2^s-1)\zeta(s)$
extends meromorphically and proves entireness. This also proves that
no unexamined higher gamma pole invalidates the converse.

The same statements in the original receiver coordinate are
\begin{equation}\tag{HTB20}
\begin{gathered}
 z_m=(4m+1)i,\qquad
 \operatorname*{Res}_{z=z_m}H_t^{\mathrm{Hur}}(z)
 =\frac{r_m(t)}{4i}
 =\frac{i(-1)^m\pi^m}{2(m-1)!}B_{2m+1}(1+t),\\
 H_t^{\mathrm{Hur}}(i)=\frac{2t+1}{16},\qquad
 H_t^{\mathrm{Hur}}(-i)=\frac1{16}.
\end{gathered}
\end{equation}
The residue factor is $1/(8s'(z_m))=1/(4i)$; it is different from
the value factor $1/8$. In particular the tangent to the Hurwitz
trajectory at zero has a genuine pole:
\begin{equation}\tag{HTB21}
 \operatorname*{Res}_{s=-2}\left.\partial_t\Xi_t(s)\right|_{t=0}
 =\pi,\qquad
 \operatorname*{Res}_{z=5i}\left.\partial_tH_t^{\mathrm{Hur}}(z)
                                      \right|_{t=0}=\frac{\pi}{4i}.
\end{equation}
This is obtained by differentiating the explicit polynomial $r_1(t)$.
In contrast $\partial_h H_h^{\mathrm{heat}}|_{h=0}$ is entire by
HTB10. Thus HTB17 is nonzero even on the actual initial zeta function;
the difference is witnessed by a specified principal part, without a
change of coordinate, multiplication factor, or zero convention.

Define the meromorphic reflection defect by
\begin{equation}\tag{HTB22}
\begin{gathered}
 R_t(s)=\Xi_t(s)-\Xi_t(1-s),\qquad R_t(1-s)=-R_t(s),\\
 R_t(0)=t,\qquad R_t(1)=-t,\qquad
 H_t^{\mathrm{Hur}}(z)-H_t^{\mathrm{Hur}}(-z)=\frac18R_t(s(z)).
\end{gathered}
\end{equation}
For real $t>-1$, its identically vanishing locus is exactly $t=0$:
necessity follows by evaluation at zero, and sufficiency follows from
HTB9. At the other entire parameter,
$R_{-1/2}(s)=(2^s-2^{1-s})\xi(s)$, which is nonzero at $s=0$.
If $r_m(t)\ne0$, $R_t$ has simple poles at both $-2m$ and $1+2m$,
with residue $r_m(t)$ at each. The reflection changes the sign of a
residue and the subtraction changes it again; the two pole sets are
disjoint. For real $t$, $\Xi_t(\bar s)=\overline{\Xi_t(s)}$ and
therefore
$H_t^{\mathrm{Hur}}(-\bar z)=\overline{H_t^{\mathrm{Hur}}(z)}$.
It is not legitimate to replace the last identity by reality on the
real $z$-axis when the reflection defect is nonzero.

The reflection is also an exact operator on the source field:
\begin{equation}\tag{HTB23}
\begin{gathered}
 (J_A f)(s)=\frac{A(1-s)}{A(s)}f(1-s),\qquad J_A^2=I,
 \qquad UJ_A=J_zU,\quad (J_zG)(z)=G(-z),\\
 [D,J_A]f(s)
 =s\frac{A(-s)}{A(s+1)}f(-s)
 -(1-s)\frac{A(1-s)}{A(s)}f(2-s).
\end{gathered}
\end{equation}
Both assertions follow by substitution; their denominators are nonzero
elements of the meromorphic field. In particular
$R_t/8$ in the receiver is $(I-J_z)UF_t$, or $U(I-J_A)F_t$.
This gives a map of the entire reflection defect, including all its
principal parts. The heat generator commutes with $J_z$, whereas its
Hurwitz counterpart has exactly the commutator transported from
HTB23. The two projections $(I\pm J_z)/2$ retain the complete even
and odd meromorphic parts; no part is discarded in their direct sum.

\subsection{The exact first transverse velocity at a critical zero}

Let $\rho$ be a simple nontrivial zeta zero with
$\operatorname{Re}\rho=1/2$. This subsection concerns each such zero;
it assumes neither that every nontrivial zero is simple nor that every
nontrivial zero lies on the line. At $\rho$, $A$ is a holomorphic unit,
so HTB14 applies. Joint holomorphy from HTB4 and the implicit function
theorem give the unique local root $\rho(t)$ with
$F_t(\rho(t))=0$, $\rho(0)=\rho$. Differentiating that equation gives
\begin{equation}\tag{HTB24}
 \rho'(0)=c_1(\rho)
 =\frac{\rho\zeta(\rho+1)}{\zeta'(\rho)},\qquad
 z_\rho(t)=-2i\left(\rho(t)-\frac12\right),\qquad
 \operatorname{Im}z_\rho'(0)=-2\operatorname{Re}c_1(\rho).
\end{equation}
The numerator is nonzero: $\rho\ne0$ and the absolutely convergent
Euler product for $\zeta(\rho+1)$ has no zero because
$\operatorname{Re}(\rho+1)>1$. This proves $c_1(\rho)\ne0$, but does
not prove that its real part is nonzero.

At a zero $\xi'(\rho)=A(\rho)\zeta'(\rho)$. Reflection and conjugation
give $\xi'(1-\rho)=-\xi'(\rho)$ and
$\overline{\xi'(\rho)}=\xi'(\bar\rho)=-\xi'(\rho)$, so this derivative
is nonzero and purely imaginary. Similarly
$c_1(1-\rho)=\overline{c_1(\rho)}$. Differentiating HTB22 at the fixed
point $\rho$ therefore yields the exact identity
\begin{equation}\tag{HTB25}
\begin{split}
 \left.\partial_tR_t(\rho)\right|_{t=0}
 &=-\xi'(\rho)c_1(\rho)
     +\xi'(1-\rho)c_1(1-\rho)\\
 &=-2\xi'(\rho)\operatorname{Re}c_1(\rho),\\
 \left.\partial_t\bigl(H_t^{\mathrm{Hur}}(z_\rho)
                 -H_t^{\mathrm{Hur}}(-z_\rho)\bigr)\right|_{t=0}
 &=-\frac14\xi'(\rho)\operatorname{Re}c_1(\rho),
 \qquad z_\rho=z(\rho).
\end{split}
\end{equation}
Thus first-order reflection failure at this zero and its first-order
transverse velocity vanish together, with the original nonzero factor
shown. This proves an exact relation at every critical simple zero;
it does not extrapolate a finite numerical nonvanishing list to all
zeros. The meromorphic global defect has already been determined
without such an extrapolation in HTB18--HTB23.

\subsection{The retained arithmetic correction and its boundary kernel}

The full difference from the original arithmetic function is a
specified entire function before completion:
\begin{equation}\tag{HTB26}
\begin{gathered}
 D_t^{\mathrm{ar}}(s)=\zeta(s)-F_t(s),\qquad
 C_t(s)=A(s)D_t^{\mathrm{ar}}(s),\qquad
 J_t(z)=\frac18C_t(s(z)),\\
 F_t+D_t^{\mathrm{ar}}=\zeta,\qquad
 \Xi_t+C_t=\xi,\qquad
 H_t^{\mathrm{Hur}}+J_t=H_0^{\mathrm{heat}}.
\end{gathered}
\end{equation}
Both summands of $D_t^{\mathrm{ar}}$ have residue $1$ at their only
possible pole, so their difference is entire. There is a convergent
boundary-kernel formula on the larger half-plane
$\operatorname{Re}s>0$:
\begin{equation}\tag{HTB27}
 D_t^{\mathrm{ar}}(s)=\frac1{\Gamma(s)}
   \int_0^{\infty}x^{s-1}d_t(x)\,dx,
 \quad
 d_t(x)=\frac{e^{-x}(1-e^{-t x})}{1-e^{-x}},\quad d_t(0)=t.
\end{equation}
At infinity both exponential terms decay for $\operatorname{Re}(1+t)>0$;
at zero Taylor expansion gives $d_t(x)=t+O(x)$. These facts prove
absolute convergence and local uniformity on the indicated domain.
On $\operatorname{Re}s>1$ it is the difference of HTB3 for $a=1$ and
$a=1+t$, so it gives the stated continuation by uniqueness. All
remaining values are fixed by the entire continuation just proved.
Differentiation gives
$\partial_tD_t^{\mathrm{ar}}=sF_t(s+1)$, also directly from HTB5.

This correction retains every completed boundary contribution:
\begin{equation}\tag{HTB28}
\begin{gathered}
 C_t(0)=-t,\qquad C_t(1)=0,\qquad
 \operatorname*{Res}_{s=-2m}C_t=-r_m(t),\\
 C_t(s)-C_t(1-s)=-R_t(s),\qquad
 \partial_t C_t(s)=sA(s)F_t(s+1).
\end{gathered}
\end{equation}
Each statement follows by subtracting the corresponding functions in
HTB26, using entireness and reflection of $\xi$ proved in HTB9, or by
differentiation. The cancellation of all the principal parts and of
the entire reflection defect is an identity of the original analytic
functions. It does not assert a zero-location theorem for any summand.

\subsection{Every map over the original lattice and its faithful receiver}

Let $L$ be any nontrivial bounded distributive lattice. There is no
finiteness, Boolean-special-case, or one-point restriction here. For a
complex vector space $V$, define the full support semimodule
\begin{equation}\tag{HTB29}
\begin{gathered}
 G_L(V)=\{(0,\lambda):\lambda\in L\}\cup(V\times\{1_L\}),\\
 (v,\lambda)+(w,\mu)=(v+w,\lambda\vee\mu),\qquad
 (c,\alpha)(v,\lambda)=(cv,\alpha\wedge\lambda),\\
 \tau=(0,0_L),\qquad e=(0,1_L).
\end{gathered}
\end{equation}
The scalars are $G_L(\mathbb C)$ with its original addition and
multiplication. Closure follows because a nonzero amplitude forces
each participating amplitude that produces it to have top support.
The semimodule axioms follow componentwise from vector-space laws and
the two distributive laws of $L$. In particular an amplitude
cancellation has the join of the original labels, rather than being
replaced by $\tau$. If $V$ is a commutative algebra, its supported
multiplication is $(v,\lambda)(w,\mu)=(vw,\lambda\wedge\mu)$; it
has exactly the same distributivity verification.

For every complex-linear map $T:V\to W$ occurring above, define
\begin{equation}\tag{HTB30}
 T^{\uparrow}(v,\lambda)=(Tv,\lambda).
\end{equation}
This is well-defined because $v=0$ implies $Tv=0$. Linearity gives
preservation of addition and of the displayed scalar action, including
all zero amplitudes. Composition and inverse maps lift exactly. Thus
$U^{\uparrow}$ is a semimodule isomorphism, the generators and their
difference commute with it as in HTB15--HTB17, and both reflection
projections and their zero-amplitude fibres lift with their original
labels, as specified in HTB33. Their semimodule kernels, meaning the
inverse images of $\tau$, are the different fibres identified there.
Where the transported algebra of HTB12 is used, $U^{\uparrow}$ is
also a unital semiring isomorphism, with unit
$(A\circ s/8,1_L)$. Its induced map on prime spectra is the usual
preimage map: an ideal is prime exactly when its inverse image is
prime, since a bijective multiplicative map preserves and reflects
membership in a product. This applies to every support prime as well
as every arithmetic prime; no prime stratum is erased by the bridge.

To verify the original faithful comparison, let $K$ be the programme's
tropical semifield with zero and Boolean subsemiring $\mathbb B$.
Its map $\beta:K\to\mathbb B$ sends zero to zero and every nonzero
element to one. It preserves products, and preserves sums because a
sum in an idempotent semifield is zero only when both summands are
zero. Regard $L$ as the $\mathbb B$-algebra with addition $\vee$ and
multiplication $\wedge$, and set
\begin{equation}\tag{HTB31}
\begin{gathered}
 K_L=K\otimes_{\mathbb B}L,\qquad
 \iota_L(\lambda)=1_K\otimes\lambda,\qquad
 \rho_L=\beta\otimes\operatorname{id}_L,\qquad
 \rho_L\iota_L=\operatorname{id}_L,\\
 \Phi_V:G_L(V)\longrightarrow V\times K_L,\qquad
 \Phi_V(v,\lambda)=(v,\iota_L(\lambda)),\\
 (T\times\operatorname{id}_{K_L})\Phi_V
       =\Phi_W T^{\uparrow}.
\end{gathered}
\end{equation}
The tensor universal property proves the retraction formula. It proves
injectivity of $\iota_L$ and therefore of $\Phi_V$. The coordinate
laws verify addition and the $G_L(\mathbb C)$-action in the product
receiver. When $V$ is an algebra they also verify multiplication.
The last line is a coordinate equality, and supplies the commuting
square for every previously specified morphism, not just for a
scalar amplitude map.

For regular germs, evaluation has its own such square. At $s_0$ where
$A$ is a unit, HTB13 lifts to
\begin{equation}\tag{HTB32}
\begin{gathered}
 \operatorname{ev}_{z_0}^{\uparrow}U^{\uparrow}(f,\lambda)
 =\left(\frac{A(s_0)}8 f(s_0),\lambda\right),\\
 \Phi_{\mathbb C}\operatorname{ev}_{z_0}^{\uparrow}U^{\uparrow}
 =\left(\operatorname{ev}_{z_0}U\times\operatorname{id}_{K_L}\right)
       \Phi_{\mathcal O_{s_0}}.
\end{gathered}
\end{equation}
At zeros or poles of $A$ the first domain is replaced by the exact
fractional module $\mathcal D_{s_0}$ of HTB13; evaluation of $Uf$ is
defined there and its lifted square still holds. The unital local
algebra arrow of HTB14 likewise gives its full $G_L$ arrow, and
retains the support of every nilpotent jet or vanishing amplitude.

In particular, at a regular zero of an original top-supported analytic
function the evaluated element is $e=(0,1_L)$, whose image is
$(0,1_{K_L})$. It is not $\tau$, whose image is $(0,0_{K_L})$.
For a linear defect operator $T$, the entire inverse image of all
zero-amplitude labels is
\begin{equation}\tag{HTB33}
 \{x:T^{\uparrow}x\in\{(0,\lambda):\lambda\in L\}\}
      =G_L(\ker T),\qquad
 (T^{\uparrow})^{-1}(\tau)=\{\tau\},\qquad
 (T^{\uparrow})^{-1}(e)=\ker T\times\{1_L\}.
\end{equation}
The first equality is just $Tv=0$ with its unchanged label. The second
uses the defining fact that label $0_L$ permits only amplitude zero;
the third uses the definition of $e$. These are three different,
fully specified fibres.

Finally the correction identity HTB26 lifts without support loss:
\begin{equation}\tag{HTB34}
 (H_t^{\mathrm{Hur}},1_L)+(J_t,1_L)
   =(H_0^{\mathrm{heat}},1_L),\qquad
 \Phi(H_t^{\mathrm{Hur}},1_L)+\Phi(J_t,1_L)
   =\Phi(H_0^{\mathrm{heat}},1_L).
\end{equation}
This uses $1_L\vee1_L=1_L$ and its tensor image; it remains true at
points where the two amplitudes cancel, provided evaluations are
taken after the regular analytic sum. At an original zero, the sum
is supported zero. Separate summands with poles are not individually
evaluated. Thus the exact analytic correction, its reflection and pole
defects, the original theta flow, and the arbitrary-lattice receiver
form one commuting construction. The construction proves the stated
maps and defects; it supplies no conclusion that all original zeta
zeros lie on, or fail to lie on, the critical line.

\subsection{The exact boundary correction to the logarithmic receiver}

For every real $t\ge-1/2$, $H_t^{\mathrm{Hur}}(0)>0$. To prove this
without a numerical input, HTB7 and HTB10 give
$H_0^{\mathrm{heat}}(0)>0$. Since
$A(1/2)=-\Gamma(1/4)/(8\pi^{1/4})<0$, it follows that
$\zeta(1/2)<0$. At $a=1/2$ the identity used in HTB19 gives
$\zeta(1/2,1/2)=(\sqrt2-1)\zeta(1/2)<0$, while HTB5 gives
$\partial_a\zeta(1/2,a)=-\zeta(3/2,a)/2<0$ for every $a>0$;
the last sign follows from its absolutely convergent positive series.
Thus $F_t(1/2)<0$ for $a\ge1/2$, and the assertion follows from HTB1.

Fix such a $t$. On a sufficiently small disc about $z=0$, both
$H_0=H_0^{\mathrm{heat}}$ and $H_t=H_t^{\mathrm{Hur}}$ are
holomorphic and nonzero. Define there $c_t=J_t/H_0$, retaining the
correction $J_t$ of HTB26. Direct differentiation gives
\begin{equation}\tag{HTB35}
\begin{gathered}
 \frac{H_t'}{H_t}-\frac{H_0'}{H_0}
     =-\frac{c_t'}{1-c_t},\qquad
 W_t(z):=\frac{H_t(z)H_t(-z)}{H_0(z)^2}
     =(1-c_t(z))(1-c_t(-z)),\\
 b_n(t):=-\frac12[z^{2n-1}]
       \left(\frac{H_t'(z)}{H_t(z)}-
             \frac{H_t'(-z)}{H_t(-z)}\right),\quad n\ge1,\\
 b_n(t)=b_n(0)-\frac12[z^{2n-1}]\frac{W_t'(z)}{W_t(z)}.
\end{gathered}
\end{equation}
The brackets mean the indicated Taylor coefficient at zero. No branch
of a global logarithm is used. For the last identity differentiate
$H_t(z)H_t(-z)=H_0(z)^2W_t(z)$, using evenness of $H_0$. This proves
the correction to every coefficient of the even logarithmic receiver,
with no finite truncation. All terms are fixed by HTB27 and the
original theta integral.

The $b_n(t)$ are defined here as logarithmic coefficients, not as
uncorrected zero-only moments of a meromorphic function. The exact
distinction is its local divisor: if
$H_t(z)=(z-z_*)^{\nu}u(z)$ with $u(z_*)\ne0$, then
$H_t'/H_t=\nu/(z-z_*)+u'/u$. A zero of multiplicity $m$ contributes
$\nu=m$ and an actual pole contributes $\nu=-1$. HTB18--HTB20 specify
all those negative contributions and every cancellation. Thus any
global zero-only interpretation must retain exactly these pole
terms. HTB35 constructs the complete local correction to the
arithmetic receiver, and HTB31 applies to its linear coefficient
extractions with every support label unchanged. It does not assert
that the nonlinear logarithm is a linear morphism of split carriers.
